\chapter{Conception de la solution}

La phase de conception est à définir très tôt pour assurer la cohérence des données et la performance du système, surtout pour la gestion des relations entre les locaux et leurs occupants.

\section{Modélisation des données (UML)}
Nous avons opté pour une modélisation centrée sur le concept de "Local", qui est le pivot de l'application. Le diagramme de classe ci-dessous détaille la structure de notre base de données PostgreSQL.

% --- Emplacement pour l'image du diagramme ---
\begin{figure}[h]
    \centering
    % \includegraphics[width=0.9\textwidth]{images/diagramme_classe.png}
    \caption{Diagramme de classe UML du système de gestion}
    \label{fig:uml_classes}
\end{figure}

\section{Description des classes principales}

\begin{description}
    \item[Local] : Cette classe recense tous les espaces physiques du département informatique. Chaque local possède un identifiant unique (ex: G206), un type spécifique dictant ses règles d'usage, et une capacité maximale d'occupation.
    \item[Personnel] : Cette entité gère à la fois les informations des membres de la faculté et leurs droits d'accès au système. L'attribut \texttt{role} permet de distinguer les administrateurs, seuls capables de modifier la liste du personnel, des utilisateurs standards.
    \item[Matériel et Prises] : Ces classes sont liées aux locaux par une relation de composition. Un équipement est physiquement rattaché à une salle, bien que le système permette son déplacement d'un local à un autre.
\end{description}

\section{Gestion des contraintes métiers}
Le diagramme prend en compte plusieurs contraintes spécifiques au sujet :
\begin{itemize}
    \item \textbf{Contrainte d'occupation} : Le nombre de personnels liés à un local ne peut excéder l'attribut \texttt{capacite\_max} du dit local.
    \item \textbf{Polyvalence des locaux} : Le type d'un local est modifiable (ex: passer d'une salle de classe à un bureau), ce qui impacte dynamiquement les équipements autorisés.
\end{itemize}